\section{问题一：确定性典型日储能线性规划}

\subsection{模型建立}

问题一的目标是在已知典型日电价、负载与光伏预测时，最小化全天外网购电费：
\[
\min C_1=\sum_{t=1}^{144}p_tg_t.
\tag{3}
\]
约束为式（1）--（2）、储能边界、购电非负，以及
\[
E_0=E_{144}=6000\ \text{kWh}.
\tag{4}
\]
该模型属于典型的微网经济调度线性规划，可直接由 HiGHS 求得全局最优解；类似的多时段微网调度结构在已有研究中也被广泛采用。\citep{umeozor2016}

\begin{figure}[H]
\centering
\includegraphics{../figures_reaslab/fig03_price_profile.pdf}
\caption{典型日外网电价曲线}
\end{figure}

\subsection{对照方案与结果}

无储能对照方案直接按 $\max(L_t-P_t,0)\Delta t$ 购电，全天费用为 48,052.047 元。优化后全天购电量为 59,482.699 kWh，购电费降为 35,126.949 元，节省 12,925.098 元，即 \textbf{26.90\%}。购电量略低于无储能对照方案的同时，主要收益来自将外网购电从高价时段转移到低价时段。

\begin{longtable}[]{@{}lrlrlr@{}}
\toprule
时间段 & 购电量/kWh & 时间段 & 购电量/kWh & 时间段 & 购电量/kWh \\
\midrule
\endhead
10:00-10:10 & 0.000 & 12:00-12:10 & 486.403 & 14:00-14:10 & 0.000 \\
16:00-16:10 & 394.931 & 18:00-18:10 & 636.983 & 20:00-20:10 & 0.000 \\
\bottomrule
\end{longtable}

\textbf{全天购电量 59,482.699 kWh；全天购电费 35,126.949 元。}

\begin{longtable}[]{@{}lrrlrr@{}}
\toprule
时间段 & 充电量/kWh & 放电量/kWh & 时间段 & 充电量/kWh & 放电量/kWh \\
\midrule
\endhead
0:00-4:00 & 4500.000 & 0.000 & 4:00-8:00 & 833.333 & 6365.841 \\
8:00-12:00 & 4787.964 & 1702.997 & 12:00-16:00 & 5286.035 & 91.101 \\
16:00-20:00 & 0.000 & 5780.132 & 20:00-24:00 & 5333.333 & 2859.868 \\
\bottomrule
\end{longtable}

\textbf{0:00 储电量 6000.000 kWh；24:00 储电量 6000.000 kWh。}

\begin{figure}[H]
\centering
\includegraphics{../figures_reaslab/fig04_q1_dispatch.pdf}
\caption{问题一储能调度对外网购电曲线的重塑}
\end{figure}

由问题一调度曲线可见，储能在低价区间和光伏富余区间充电，在高价/高净负荷时段放电；SOC 始终位于允许区间内，并在日末回到 6000 kWh，符合题设。
